% =====================================================================
%  The Evidence Peer-Review System
%  An engineering paper for The Disclosure Platform. Specifies the
%  smart contract (EvidenceConsensus on BNB Smart Chain mainnet), its
%  read-only Lens sidecar (EvidenceConsensusLens) and its peer-membership
%  sidecar (PeerGovernance), the off-chain projection (Supabase),
%  content/taxonomy hashing, the binding model, the eight-state
%  lifecycle, the bounded submission queue, the threshold mathematics,
%  order-independence via peer-count snapshots, the challenge flow, the
%  administrator-free peer registry, the capture boundary and
%  force-renounce backstop, peer-liveness garbage collection, taxonomy
%  retirement, and the universal EIP-712 vote-by-signature scheme that
%  authorises every peer act on-chain (with a hashed off-chain
%  deliberation note) --- and shows how each delivers the goals set out
%  in the companion essay "Disclosure: A Labour of Love".
%  Class : article
% =====================================================================

\documentclass[11pt,a4paper]{article}

% ---------- Packages ----------
\usepackage[T1]{fontenc}
\usepackage[utf8]{inputenc}
\usepackage{lmodern}
\usepackage[a4paper,margin=1in]{geometry}
\usepackage{microtype}
\usepackage{parskip}
\usepackage{titlesec}
\usepackage{enumitem}
\usepackage{booktabs}
\usepackage{array}
\usepackage{xcolor}
\usepackage{textcomp}
\usepackage[skins,breakable]{tcolorbox}
\usepackage{amsmath}
\usepackage{amssymb}
\usepackage{graphicx}
\usepackage[hidelinks,colorlinks=true,linkcolor=black,urlcolor=blue!55!black,citecolor=black]{hyperref}
\usepackage{fancyhdr}

% ---------- Colours ----------
\definecolor{ink}{HTML}{0E1116}
\definecolor{inkfaint}{HTML}{555B66}
\definecolor{accent}{HTML}{6A4CFF}
\definecolor{accent2}{HTML}{1F8F66}
\definecolor{warn}{HTML}{B85C00}
\definecolor{rule}{HTML}{C8CCD3}
\definecolor{panelbg}{HTML}{F6F4EE}

% ---------- Helper macros ----------
\newcommand{\arr}{\ensuremath{\rightarrow}}
\newcommand{\ceilop}[1]{\ensuremath{\left\lceil #1 \right\rceil}}
\newcommand{\floorop}[1]{\ensuremath{\left\lfloor #1 \right\rfloor}}

% Let TeX relax a line as a last resort so long monospace identifiers in
% prose (e.g. \texttt{motionForceRenounce}) never spill into the margin.
\setlength{\emergencystretch}{2em}

% ---------- Page style ----------
\pagestyle{fancy}
\fancyhf{}
\fancyhead[L]{\footnotesize\textsc{The Disclosure Platform}}
\fancyhead[R]{\footnotesize\textsc{The Evidence Peer-Review System}}
\fancyfoot[C]{\footnotesize\thepage}
\renewcommand{\headrulewidth}{0.3pt}
\renewcommand{\footrulewidth}{0pt}

% ---------- Section style ----------
\titleformat{\section}
  {\normalfont\Large\bfseries\color{ink}}
  {\thesection}{0.6em}{}
\titleformat{\subsection}
  {\normalfont\large\bfseries\color{ink}}
  {\thesubsection}{0.6em}{}

% ---------- Helper boxes ----------
\newtcolorbox{plainbox}[1][]{
  enhanced, breakable, boxrule=0pt, frame hidden,
  colback=panelbg, colupper=ink,
  left=10pt, right=10pt, top=8pt, bottom=8pt,
  arc=2pt, #1
}

\newtcolorbox{notebox}[1][]{
  enhanced, breakable, boxrule=0pt, frame hidden,
  colback=panelbg, colupper=ink,
  left=12pt, right=12pt, top=8pt, bottom=8pt,
  borderline west={2pt}{0pt}{accent},
  arc=0pt, #1
}

% ---------- Title block ----------
\title{
  \vspace{-2em}
  {\Huge\bfseries The Evidence Peer-Review System}\\[0.6em]
  {\Large Engineering an open, tamper-evident,\\
   capture-resistant, evolving consensus}
}
\author{\textbf{The Disclosure Platform}}
\date{May 2026}

% =====================================================================
\begin{document}
% =====================================================================
\maketitle
\thispagestyle{empty}

\vspace{-1em}
\begin{plainbox}
\textbf{Abstract.}
This is an engineering paper. It specifies the system that turns the
philosophy of \emph{Disclosure: A Labour of Love} into enforceable
mechanism. A smart contract, \texttt{EvidenceConsensus}, deployed on
BNB Smart Chain \emph{mainnet}, is the sole source of truth for the peer
set, the Pillar~\arr~Topic taxonomy, and the lifecycle of every record;
a read-only sidecar, \texttt{EvidenceConsensusLens}, serves aggregate
views without holding any authority, and a second sidecar,
\texttt{PeerGovernance}, holds the nominee/revocation flows so the core
stays under the EIP-170 deploy ceiling. A Supabase database is a fast,
rebuildable \emph{projection} of the chain, kept in step by an indexer
that joins rows to on-chain events by deterministic hashes. We specify
how a record's content is hashed so any later tampering is detectable;
how the taxonomy is grown by peer consensus and why a node is never
empty; the \emph{binding} --- an (evidence~$\times$~topic) pair --- as
the atomic unit of review; the eight-state lifecycle a binding moves
through; the bounded submission queue that keeps open submission from
outrunning review capacity; the threshold mathematics that scales with
the active peer set and makes ``Canon'' a true majority; how vote
outcomes are made order-independent via peer-count snapshots; the
time-boxed challenge procedure; the administrator-free peer registry;
the BFT-style capture boundary and the peer \emph{force-renounce}
backstop against a captured owner; peer-liveness garbage collection;
taxonomy retirement; and the \emph{universal} EIP-712
vote-by-signature scheme --- a single \texttt{Vote} typehash on the core
and a single \texttt{PeerVote} typehash on the governance sidecar ---
that authorises \emph{every} peer act on-chain, each carrying an
optional, signed off-chain deliberation note. We close by mapping each
mechanism back to the principle it secures.
\end{plainbox}

\tableofcontents
\vspace{1em}
\hrule
\vspace{1em}

% =====================================================================
\section{Architecture: a chain of truth and a cache of speed}
% =====================================================================

The system is two layers with a strict division of authority, plus two
sidecars at the chain layer that hold either no authority at all (the
Lens) or only a narrowly scoped, peer-gated authority over the peer
registry (the governance sidecar).

\textbf{Layer 1 --- the chain of truth.} A Solidity core,
\texttt{EvidenceConsensus}, on BNB Smart Chain (BSC) mainnet, holds
everything that must be authoritative and tamper-resistant: the set of
verified peers, the taxonomy of pillars and topics, the state and vote
tallies of every record, and the time windows that govern them. Nothing
is canon unless the chain says so. Because the chain is public and
append-only, the entire deliberative history is preserved end to end and
cannot be silently rewritten.

\textbf{The read sidecar.} As the core grew --- a submission queue and
on-chain peer garbage collection were added --- its deployed bytecode
approached the EIP-170 limit of 24{,}576 bytes. Pure aggregation views
(assembling the peer list, the nominee list, and the pending-proposal list
with their tallies) were therefore moved into a separate contract,
\texttt{EvidenceConsensusLens}. The Lens holds \emph{no storage}, has
\emph{no privileges}, and reconstructs every aggregate purely from the
core's public state, so it carries zero consensus risk: the core remains
the sole writer and the sole source of truth. Read clients call the Lens
for these list views.

\textbf{The peer-governance sidecar.} When universal EIP-712
vote-by-signature was added to the hot voting path (Section~14), the
core ran out of EIP-170 headroom again. The nominee admission and
revocation flows --- which do not touch the binding vote path; they only
need the core's peer \emph{registry} state --- were moved into a
separate contract, \texttt{PeerGovernance}. The sidecar holds its own
nominee/revocation bookkeeping and its own EIP-712
\texttt{PeerVote} domain, and mutates the core's peer set through two
privileged hooks, \texttt{gAddPeer} / \texttt{gRemovePeer}, gated by an
\texttt{onlyGovernance} modifier wired once at deploy via
\texttt{setGovernance}. The core therefore remains the source of truth
for who is a peer; the sidecar is the only authorised path (besides
seed-phase owner \texttt{addPeer} and inactivity GC) to change that set
once the seed phase is over. Read clients hit the Lens for list views,
the sidecar for membership writes, and the core for everything else
that mutates state.

\textbf{Layer 2 --- the cache of speed.} A Supabase (Postgres) database
holds a \emph{projection} of the chain: the same records and taxonomy,
shaped for fast reading, searching, and pagination. The cache is never the
authority. It is reconciled to the chain by an indexer edge function,
\texttt{chain-indexer-evidence}, which reads chain events and joins them to
off-chain rows by deterministic hashes (a binding's row is matched to its
on-chain event by \texttt{binding\_hash}; a taxonomy row by
\texttt{node\_hash}). The whole database is built from a consolidated base
migration plus a few additive follow-ons, with \emph{zero seed data} --- it can
be rebuilt from an empty Postgres instance and re-synchronised from the chain.

\textbf{The frontend.} A Vite multi-page React application reads from the
cache for speed and writes to the contract through the reader's own wallet
(MetaMask). The heavy chain library (\texttt{ethers}~v6) is loaded lazily,
so a visitor who only reads the archive never pays its cost. Reads can
optionally use a configured read RPC; writes always go through the user's
wallet, so the user --- never the platform --- signs and pays for every state
change they cause.

\begin{notebox}
\textbf{Why this split.} Putting authority on-chain makes the record
tamper-evident and ownerless; putting presentation off-chain makes it fast
and pleasant; splitting list aggregation into a privilege-free Lens and
narrow peer-registry mutation into a peer-gated governance sidecar keeps
the core small enough to deploy without weakening it. Keeping the cache
strictly subordinate --- rebuildable, never authoritative --- means a
compromised or lost database costs convenience, not truth. The truth is
whatever the contract can prove.
\end{notebox}

% =====================================================================
\section{The record and content hashing}
% =====================================================================

A piece of evidence has one canonical identity and one content fingerprint.
The identity is a UUID, mapped one-to-one to a \texttt{bytes32} for
on-chain use. The fingerprint is a keccak-256 hash over the record's
\emph{content fields only}:

\begin{plainbox}
\texttt{contentHash = keccak256( utf8( canonicalJSON \{}\\
\hspace*{2em}\texttt{title, source, year, excerpt, link, tier}\\
\texttt{\} ) )}
\end{plainbox}

The canonical JSON trims each string field and coerces \texttt{tier} to a
number, so the hash is reproducible byte-for-byte by any party. Two
properties of this design matter.

\textbf{The topic is deliberately excluded from the hash.} A record's
fingerprint binds \emph{what the evidence is}, not \emph{where it is filed}.
This is what lets the same evidence be cross-listed under many topics
without ever being re-hashed --- one piece of evidence, one
\texttt{contentHash}, many filings.

\textbf{The hash is computed identically in three places.} The frontend
computes it when a record is submitted; the \texttt{verify-attestation} edge
function recomputes it when a vote is verified; and the
\texttt{audit-content-hash} edge function recomputes it on a schedule across
the archive. All three must produce identical bytes. If the content stored
off-chain is ever altered, its recomputed hash no longer matches the hash
committed on-chain, and the discrepancy is recorded as a tamper alert.
Tampering is not prevented by trust; it is made \emph{detectable} by
arithmetic.

% =====================================================================
\section{The Pillar--Topic taxonomy}
% =====================================================================

Evidence is not filed into a flat list. It is filed into a two-level
taxonomy that the peers themselves grow: \textbf{pillars} (top-level
domains) and, beneath each pillar, \textbf{topics}. The archive grows
\emph{wider} as pillars are added and \emph{deeper} as topics are added, and
both kinds of growth happen only by peer consensus.

Every taxonomy node has a deterministic on-chain identity and an off-chain
metadata commitment:

\begin{plainbox}
\texttt{nodeId   = keccak256( utf8( slug ) )}\\[0.3em]
\texttt{metaHash = keccak256( utf8( canonicalJSON \{}\\
\hspace*{2em}\texttt{kind, slug, parent, title, blurb, tag}\\
\texttt{\} ) )}
\end{plainbox}

Because the id is the hash of the slug, the indexer can join an off-chain
taxonomy row to its on-chain node by \texttt{node\_hash} without any shared
database key. Both commitments are tamper-audited exactly as evidence content
is (Section~2): the \texttt{audit-content-hash} job recomputes every
pillar/topic's \texttt{node\_hash} and \texttt{metaHash} on a schedule and
raises a tamper alert if the off-chain row has drifted from what the peers
committed on-chain, so a silently edited taxonomy node is as detectable as a
silently edited record.

\subsection{A taxonomy node is never empty}

A central rule prevents the archive from filling with empty scaffolding:
\emph{a node is never proposed alone.}

\begin{itemize}[leftmargin=1.4em]
  \item A \textbf{pillar} is proposed bundled with its first \textbf{topic}
  \emph{and} a founding piece of \textbf{evidence}.
  \item A \textbf{topic} is proposed bundled with a founding piece of
  \textbf{evidence}.
\end{itemize}

The whole bundle rides on a single endorsement gate. The relevant contract
calls carry the entire bundle at once:

\begin{plainbox}
\texttt{proposePillar(id, metaHash, topicId, topicMetaHash,}\\
\hspace*{2em}\texttt{evidenceId, tier, contentHash)}\\[0.3em]
\texttt{proposeTopic(id, parentPillar, metaHash,}\\
\hspace*{2em}\texttt{evidenceId, tier, contentHash)}
\end{plainbox}

A proposal is then endorsed by other peers via \texttt{endorseNode(id)} (the
proposer counts as endorsement number one). When endorsements reach the
bundle threshold (Section~7), the node ratifies. For a pillar, its bundled
founding topic ratifies in the same act, and while the pillar is still
pending it \emph{reserves} both its bundled topic id and its bundled
evidence id so nothing else can claim them. Critically, the founding
evidence's (evidence~$\times$~topic) binding is canonized \emph{atomically}
with ratification --- it goes straight to Canon, because the taxonomy
endorsement \emph{is} the consensus for that founding record.

\begin{notebox}
\textbf{Bundling is not a cheaper way in.} The founding evidence is not
exempt from review-grade consensus. The bundle gate is
$\max(\text{taxonomyThreshold},\ \text{canonizeThreshold(tier)})$, so
smuggling a record in as a thin ``new topic'' is never cheaper than the
normal review path. A pillar with no topic, or a topic with no evidence, is
an opinion about how the archive \emph{should} be organised with nothing yet
to organise. Bundling forces every act of structure to arrive already
carrying content, so the taxonomy can only grow where there is real evidence
to fill it. Further evidence added later to a ratified topic
(\texttt{submitEvidence} / \texttt{fileBinding}) passes through ordinary
tier-based review.
\end{notebox}

A proposal that never reaches its gate is not stranded forever: after the
30-day \texttt{PROPOSAL\_WINDOW} anyone may call \texttt{lapseProposal(id)},
which garbage-collects the node and frees its reserved topic and evidence
ids for a fresh attempt. Each peer is also capped at a small number of
outstanding (still-proposed) proposals, the on-chain analogue of the
off-chain \texttt{tax\_insert:} throttle, so no single peer can flood the
proposal queue.

% =====================================================================
\section{Bindings: the atomic unit of review}
% =====================================================================

The system does not review ``evidence'' in the abstract. It reviews a
\textbf{binding}: a specific (evidence~$\times$~topic) pair. Each binding is
an independent voting unit with its own lifecycle state and its own tallies,
and its own deterministic id:

\begin{plainbox}
\texttt{bindingId = keccak256( abi.encode( evidenceId, topicId ) )}
\end{plainbox}

This is the structural consequence of ``one record, many filings.'' A single
piece of evidence (one \texttt{contentHash}) may be filed under any number of
topics, and each filing is judged on its own merits. The same record can be
Canon under one topic and Expelled under another, because relevance is
contextual: a document may be excellent evidence for one subject and
irrelevant to another.

Two contract calls create bindings. \texttt{submitEvidence(id, tier,
topicId, contentHash)} registers a new record and opens its first binding.
\texttt{fileBinding(id, topicId)} cross-lists an already-registered record
under a further ratified topic, opening a new independent binding. Every
subsequent vote or challenge call is addressed to a binding by its
\texttt{(id, topicId)} pair.

\begin{notebox}
\textbf{Open submission.} \texttt{submitEvidence} and \texttt{fileBinding}
are callable by any wallet; filing evidence is not gated. What is gated is
\emph{judgement} --- only verified peers can vote a binding through its
lifecycle. This is the contract-level expression of ``anyone may file; named
peers verify.''
\end{notebox}

% =====================================================================
\section{The eight-state lifecycle}
% =====================================================================

A binding moves through a state machine. The contract enumerates nine
values; \texttt{None} is merely the default for a binding that does not
exist, leaving eight reachable states.

\begin{center}
\begin{tabular}{@{}rll@{}}
\toprule
\# & \textbf{State} & \textbf{Meaning} \\
\midrule
0 & None       & The binding does not exist (default). \\
1 & Submitted  & In active review; the pending window is open. \\
2 & Canon      & Canonized by peer consensus; in the public archive. \\
3 & Expelled   & Rejected by peer consensus (terminal). \\
4 & Lapsed     & Review window passed without a threshold; re-filable. \\
5 & Contested  & A canon binding under an active challenge. \\
6 & Deprecated & Removed by challenge consensus (terminal). \\
7 & Reaffirmed & Survived a challenge; re-confirmed as canon. \\
8 & Queued     & Parked behind a full active-review set; not yet in review. \\
\bottomrule
\end{tabular}
\end{center}

\vspace{0.6em}

The legal transitions are:

\begin{plainbox}
\centering
\texttt{Queued} \;\arr\; \texttt{Submitted}\\[0.4em]
\texttt{Submitted} \;\arr\; \texttt{Canon} \;$\mid$\; \texttt{Expelled} \;$\mid$\; \texttt{Lapsed}\\[0.4em]
\texttt{Lapsed} \;\arr\; \texttt{Submitted/Queued} \;\;(\text{re-filed})\\[0.4em]
\texttt{Canon\,/\,Reaffirmed} \;\arr\; \texttt{Contested} \;\arr\; \texttt{Deprecated} \;$\mid$\; \texttt{Reaffirmed}
\end{plainbox}

A newly opened binding enters \texttt{Submitted} if the active-review set has
a free slot, otherwise it parks in \texttt{Queued} until promoted
(Section~6). In \texttt{Submitted} it sits for a 30-day review window.
Approvals reaching the canonize threshold move it to \texttt{Canon};
rejections that make canon arithmetically impossible settle it early; if
neither happens before the window closes, anyone may call \texttt{markLapsed}
to settle it. The early-settle and window-close paths share one rule
(Section~8): an expel-quorum of rejections yields \texttt{Expelled}
(terminal), otherwise \texttt{Lapsed} (re-filable with a fresh review round).
Only \texttt{Canon} and \texttt{Reaffirmed} bindings appear in the public
archive. A canon binding may later be challenged into \texttt{Contested},
resolving to \texttt{Deprecated} or \texttt{Reaffirmed} (Section~9).
Expelled, Lapsed, and Deprecated bindings are not erased --- they remain
on-chain as part of the permanent record of what the network decided.

\begin{center}
\begin{tabular}{@{}lll@{}}
\toprule
\textbf{Window / cooldown (constant)} & \textbf{Scope} & \textbf{Duration} \\
\midrule
\texttt{PENDING\_WINDOW} (review)        & per binding & 30 days \\
\texttt{CHALLENGE\_WINDOW}               & per binding & 21 days \\
\texttt{CHALLENGE\_COOLDOWN}             & per peer    & 7 days  \\
\texttt{RECHALLENGE\_COOLDOWN}           & per binding & 30 days \\
\texttt{PROPOSAL\_WINDOW} (taxonomy)     & per node    & 30 days \\
\texttt{REVOKE\_WINDOW}                  & per motion  & 14 days \\
\texttt{INACTIVITY\_WINDOW} (peer GC)    & per peer    & 30 days \\
\texttt{SUBMIT\_COOLDOWN} (non-peer)     & per address & 10 minutes \\
\texttt{BOOST\_COOLDOWN} (non-peer)      & per address & 10 minutes \\
\bottomrule
\end{tabular}
\end{center}

% =====================================================================
\section{The submission queue: open intake, bounded review}
% =====================================================================

Open submission scales without limit; human review does not. If every
submission entered active review immediately, a burst of filings --- honest
or hostile --- could bury the peers under more open reviews than they could
ever process, and the 30-day clock on each would start ticking regardless.
The contract bounds the \emph{active} review set instead, as a fixed
shared batch the whole network works through in the same order.

\begin{plainbox}
\texttt{reviewCapacity() = REVIEW\_CAPACITY}
\hspace{1em}(\texttt{REVIEW\_CAPACITY = 3})
\end{plainbox}

When a binding is opened, it enters \texttt{Submitted} only if
\texttt{activeReviewCount < reviewCapacity()}; otherwise it parks in
\texttt{Queued} with its review clock unset. The shared-batch design is
deliberate: every peer is looking at the same small set of bindings, so
the network converges on each verdict together rather than spreading
attention thin across hundreds of half-reviewed records. As bindings
resolve (canonized / expelled / lapsed) they free slots, and queued
bindings are promoted to fill them:

\begin{itemize}[leftmargin=1.4em]
  \item \texttt{boostQueued(id, topicId)} casts \emph{one} boost on a queued
  binding to raise its priority. It is open to \emph{anyone} --- a boost only
  reorders the queue (the keeper promotes by priority) and never touches the
  consensus verdict, so it is open triage, not a vote. Two anti-spam guards
  bound it: one boost per wallet per binding, and a per-wallet
  \texttt{BOOST\_COOLDOWN} (10 minutes) so a single identity cannot sweep-boost
  the whole queue; active peers are exempt from the cooldown, mirroring the
  public submit policy. Because a boost can only change the \emph{order} in
  which already-queued evidence is reviewed --- never what canonizes --- the
  worst a sybil flood achieves is reordering the queue, which the permissionless
  drain (below) makes self-correcting.
  \item \texttt{promote(id, topicId)} moves the highest-priority queued
  binding into \texttt{Submitted} once a slot is free, starting its review
  clock. Promotion is \emph{permissionless}: a keeper (the
  \texttt{consensus-keeper} edge function, Section~15) calls it in priority
  order, but anyone may call it, so the queue always drains even if the
  keeper is offline. Only ordering --- never liveness --- depends on the
  keeper.
\end{itemize}

This global bound, together with a per-address \texttt{SUBMIT\_COOLDOWN} of
10 minutes on non-peer submitters, replaces an earlier per-submitter
outstanding cap: a single identity can no longer flood the \emph{active}
set no matter how much it queues, and the active review load the network
carries at any moment is fixed by the contract rather than by the size
of the inbox.

\begin{notebox}
\textbf{Self-healing by construction.} \texttt{activeReviewCount} falls every
time a binding leaves \texttt{Submitted}, which immediately frees a slot for
the next promotion. The queue therefore drains at exactly the rate the peers
clear reviews, and it cannot deadlock: every path out of \texttt{Submitted}
decrements the counter, and \texttt{promote} is callable by anyone.
\end{notebox}

% =====================================================================
\section{Voting and threshold mathematics}
% =====================================================================

Every threshold is a function of \texttt{activePeerCount}, the live size of
the verified peer set, so the bar tracks the network as it grows and the
system is never deadlocked when it is small. Every threshold carries a
\textbf{floor of 1}, so a single Genesis peer can act without stalling.

For canonize, expel, and deprecate, the threshold is the ceiling of a
percentage of the peer count:
\[
  t(n) \;=\; \max\!\left(1,\; \ceilop{\,n \cdot p\,}\right),
\]
where $n$ is the relevant peer count and $p$ is the percentage below.

\begin{center}
\begin{tabular}{@{}lccc@{}}
\toprule
\textbf{Action} & \textbf{Tier I} & \textbf{Tier II} & \textbf{Tier III} \\
\midrule
Canonize (approvals to accept)   & 60\% & 55\% & 51\% \\
Expel (rejections to reject)     & \multicolumn{3}{c}{25\% (all tiers)} \\
Deprecate (challenge votes)      & 65\% & 60\% & 55\% \\
\bottomrule
\end{tabular}
\end{center}

\vspace{0.4em}

\textbf{Canon is a true majority.} Canonization clears more than half of the
peer set at every tier (60 / 55 / 51\%), so ``Canon'' reflects majority
consensus, never a large minority. The asymmetry across tiers is
intentional. The strongest evidence (Tier~I: peer-reviewed and declassified)
is the \emph{hardest} both to admit (60\%) and to remove once admitted
(65\%): the record demands a high bar to take its strongest claims seriously,
and an even higher bar to retire them. Testimony (Tier~III) is the easiest to
admit (51\%) but still requires a substantial majority (55\%) to deprecate,
so a contested first-person account is not casually discarded. Early
expulsion of a still-pending record is tier-independent and deliberately low
(25\%): a record the network is clearly unwilling to accept should not have
to wait out the full window.

Four further thresholds govern membership and structure. They are
deliberately \emph{decoupled} so that admitting a peer, ratifying a node,
removing a peer, and retiring a node each clear an appropriate and distinct
bar:

\begin{center}
\begin{tabular}{@{}lll@{}}
\toprule
\textbf{Gate} & \textbf{Formula} & \textbf{Meaning} \\
\midrule
Admission (nominee)   & $\floorop{n/3} + 1$   & strictly more than $1/3$ \\
Taxonomy ratify       & $\floorop{n/2} + 1$   & strict majority \\
Revoke a peer         & $\ceilop{n/2}$        & simple majority \\
Retire a node         & $\ceilop{2n/3}$       & $2/3$ supermajority \\
\bottomrule
\end{tabular}
\end{center}

\vspace{0.4em}

The founding-bundle gate composes two of these:
$\text{bundleThreshold(tier)} = \max\!\big(\floorop{n/2}+1,\;
\text{canonizeThreshold(tier)}\big)$, so a founding record is never
canonized on a cheaper vote than ordinary review would demand.

Double-voting is impossible: the contract records, per binding and per phase
(review or challenge), whether a given address has already voted, and each
re-filing or re-challenge bumps a round counter that isolates the new tally
from the old. Peers may also cast review votes in batches (up to
\texttt{MAX\_REVIEW\_BATCH = 50} per transaction) to process a queue
economically.

% =====================================================================
\section{Order-independence via peer-count snapshots}
% =====================================================================

A subtle hazard in any threshold-on-a-percentage scheme is that the
denominator moves: peers join and leave \emph{during} a review window. If the
verdict depended on when each vote happened relative to that churn, the same
set of votes could canonize or not by accident of timing. The contract
removes this dependence deliberately.

When a binding enters review (at submission or at promotion) it
\emph{snapshots} \texttt{activePeerCount} into \texttt{peerSnapshot}. The
\textbf{canonize target} is then judged against that frozen snapshot, so a
fixed number of approvals canonizes regardless of membership changes after
the window opened. The same snapshot governs the \textbf{expel} bar.

Early settlement, however, asks a different question, and so it uses the
\emph{live} count. A reject vote triggers settlement only once canon has
become arithmetically impossible --- when, even if every peer who has not yet
voted approved, approvals still could not reach the canonize target:
\[
  \text{canonizeTarget} + \text{rejects} \;>\; \texttt{activePeerCount}.
\]
Judging ``canon impossible'' against the live electorate keeps it mutually
exclusive with canonization under peer-set \emph{growth}: late-joining peers
can still rescue a record that a snapshot-based test would have prematurely
declared dead. When settlement does fire, an expel-quorum of rejections
($\geq$ the snapshot's 25\% bar) yields \texttt{Expelled}; otherwise the
binding merely failed to attract a verdict and \texttt{Lapses}, re-filable
later. The one residual effect of churn is benign: a heavy peer-set
\emph{shrink} can turn what would have been an expulsion into a lapse, so
membership loss can never \emph{terminally} expel otherwise-canonizable
evidence --- it can only send it back to the queue against the smaller
electorate.

\begin{notebox}
\textbf{The guarantee.} Within a review round the outcome is independent of
the order in which votes arrive: a binding canonizes the instant approvals
reach the (snapshot) target, and is expelled early only once canon is
genuinely impossible against the live set. Snapshot for the target, live
count for ``impossible'' --- the two are chosen so the two conclusions can
never both be true at once.
\end{notebox}

% =====================================================================
\section{Challenges: keeping the record revisable}
% =====================================================================

Canon is never permanent by default. Any verified peer may challenge a
\texttt{Canon} or \texttt{Reaffirmed} binding with \texttt{openChallenge(id,
topicId)}, which moves it to \texttt{Contested} and starts the 21-day
challenge window. The challenger's own vote is cast automatically as the
first challenge vote. Two cooldowns keep the mechanism from becoming a weapon:
a \emph{per-peer} \texttt{CHALLENGE\_COOLDOWN} (7 days) limits how often one
peer can open challenges, and a \emph{per-binding}
\texttt{RECHALLENGE\_COOLDOWN} means a record cannot be re-contested until its
last challenge window has closed plus 30 days --- so a rotating cast of peers
cannot keep evidence perpetually in limbo.

During the window, peers vote with \texttt{castChallengeVote(id, topicId,
support)} --- either supporting the challenge or defending the record. The
resolution is a race between a threshold and a clock:

\begin{itemize}[leftmargin=1.4em]
  \item If support votes reach the tier-dependent deprecate threshold, the
  binding becomes \texttt{Deprecated} immediately.
  \item If the 21-day window closes without that threshold being met, the
  binding is \texttt{Reaffirmed} --- it survived, and is now a record that has
  been tested and held.
\end{itemize}

Either outcome is settled with \texttt{finalizeChallenge(id, topicId)}.
Crucially, deprecation is judged against the \emph{live} peer set, not a
snapshot: destroying canon must always clear a supermajority of the network
\emph{as it stands now}, so peers admitted during the challenge window raise
the bar, never dilute it. Because a reaffirmed binding can itself be
challenged again later (after its cooldown), no record is ever beyond
re-examination; it is only ever ``standing for now.''

% =====================================================================
\section{The peer registry: membership without an administrator}
% =====================================================================

The set of peers who may judge evidence is governed entirely on-chain,
by the peers themselves. All membership writes live on the
\texttt{PeerGovernance} sidecar (Section~1); the core exposes the peer
set as a view and accepts membership changes only through the sidecar's
peer-gated \texttt{gAddPeer} / \texttt{gRemovePeer} hooks.

\textbf{Genesis and seed phase.} The system launches from a Genesis peer
set supplied at deployment (the quorum floor of 1 lets even a one-peer
network act). During an initial seed phase the owner may seed further
peers directly with the core's \texttt{addPeer}, but \emph{only} while
\texttt{activePeerCount < seedPhaseK}; the moment the active set reaches
\texttt{seedPhaseK}, ordinary nominations open
(\texttt{nominationsOpen()} on the sidecar) and the owner's seeding
power ends. If the owner has \emph{renounced}, nominations stay open
unconditionally --- otherwise a later shrink below the seed floor with
no owner to re-seed would permanently strand the network.

\textbf{Steady state.} A newcomer is brought in by an existing peer
via the sidecar's \texttt{nominatePeer} (handle, note hash, signature;
the handle is at most 64 bytes). The nominator's signature is an
EIP-712 \texttt{PeerVote} of kind \texttt{nominate} (Section~14) that
the contract recovers and matches against \texttt{msg.sender}, so a
nomination is a signed, accountable act --- not just a function call.
Other peers then endorse the nominee with \texttt{endorseNominee} (a
\texttt{PeerVote} of kind \texttt{nominee-endorse}); when endorsements
reach the admission gate $\floorop{n/3}+1$, the sidecar calls back
into the core's \texttt{gAddPeer} hook and the nominee is verified
\emph{automatically} --- no administrator approves the final step. A
nomination that never reaches its gate can be cleared after the
proposal window by the permissionless \texttt{lapseNominee}, freeing
the address to be nominated again with a fresh endorsement round.

\textbf{Removal.} A peer can be removed by the same community that
admits them: \texttt{motionRevoke} opens a motion (the motioner is
discard vote one and signs a \texttt{PeerVote} of kind \texttt{revoke}
for the new round) and \texttt{voteRevoke} accrues further signed
discards, with removal at a majority $\ceilop{n/2}$. A stale motion
that never passed is cleared after \texttt{REVOKE\_WINDOW} by
\texttt{cancelStaleRevocation}, so no peer is left indefinitely under
an open cloud. A peer-floor invariant forbids removing the last peer.

\textbf{Ownership.} The owner role exists only for the seed phase and
for emergency \texttt{pause}/\texttt{unpause}; it can never force a
record in, force a peer out, or override a vote. Ownership transfers in
a safe two-step \texttt{proposeOwner}~\arr~\texttt{acceptOwnership}
handshake (cancellable), and once the seed phase is complete and the
contract is live the owner may \texttt{renounceOwnership} to drop the
role entirely.

\begin{notebox}
\textbf{Identity is a wallet plus a handle, and now a signed reason.}
A peer is a wallet address carrying a human-readable handle. The system
makes no claim that a wallet is a verified legal identity --- but every
judgement is permanently tied to the wallet that signed it, and (since
the universal vote-by-signature rollout) every membership act also
carries an optional deliberation note hashed into that signature.
Identity here is wallet + handle + a public history of signed reasons.
\end{notebox}

% =====================================================================
\section{Capture resistance and the force-renounce backstop}
% =====================================================================

An ownerless system is only as good as its resistance to a coalition quietly
\emph{becoming} the owner. The registry is designed around an explicit
capture boundary.

\textbf{The admission gate is the capture line.} Admission requires
\emph{strictly} more than one third of peers ($\floorop{n/3}+1$). The strict
inequality is load-bearing: with a mere $\ceilop{n/3}$ a coalition of exactly
$n/3$ (when $n$ is divisible by 3) could admit its own members and bootstrap
itself to a majority. The strict gate closes that zero-margin case, so any
coalition holding $\leq 1/3$ of the peers cannot admit anyone --- nor ratify
or retire a node --- without honest cooperation. This is the familiar BFT
boundary: the system is safe as long as more than two thirds of peers are
honest.

\textbf{Creating and destroying canon taxonomy are both kept above easy
reach.} Ratifying a node is a strict majority ($\floorop{n/2}+1$) and
retiring one is a $2/3$ supermajority ($\ceilop{2n/3}$). Decoupling the
ratify gate from the lower admission gate means a sub-majority faction can no
longer spawn nodes that a supermajority must then clean up; the
create/retire gap narrows from $1/3 \to 2/3$ to $1/2 \to 2/3$, and the retire
bar sits safely above the capture line.

\textbf{The force-renounce backstop.} The owner's only powers are pause
and seed-phase \texttt{addPeer}, and \texttt{renounceOwnership} is gated
on the contract not being paused --- so a malicious or compromised owner
could, in principle, pause forever and brick the network with no peer
recourse. One motion closes that gap.
\texttt{motionForceRenounce(noteHash, sig)} and
\texttt{voteForceRenounce(noteHash, sig)} let a $2/3$ peer supermajority
(the same bar used to destroy canon governance) strip the owner
entirely \emph{and} unpause the contract. Each vote is an EIP-712
\texttt{Vote} (phase~4, with a fixed sentinel \texttt{bindingId =
keccak256("force-renounce")}) that the contract recovers on-chain and
matches against \texttt{msg.sender}, so the most consequential act in
the system is also one of its most carefully attributed --- signed,
counted, and bound to a deliberation note. These two functions
deliberately carry \emph{no} ``when-not-paused'' guard --- their whole
purpose is to act precisely when a captured owner has paused
everything. When they pass, ownership passes to no one: \texttt{owner}
is set to the zero address and the network continues fully
self-governed.

\begin{notebox}
\textbf{What this buys.} ``Owned by no one'' is not merely a promise to
renounce. Even an owner who refuses to renounce, and who pauses the contract
to entrench, can be evicted by two thirds of the peers --- the same honest
majority the whole BFT model already assumes. The keys cannot be held against
the network's will.
\end{notebox}

% =====================================================================
\section{Peer liveness and garbage collection}
% =====================================================================

Every threshold divides by \texttt{activePeerCount}. If that denominator
counted peers who had silently abandoned the network, the bars would drift
out of reach and a live, willing minority could be unable to reach any
quorum. The contract keeps the denominator honest with an objective liveness
clock.

Each peer carries a \texttt{lastActive} timestamp, set on admission and
refreshed by every consensus vote (review or challenge) and by an explicit
\texttt{heartbeat()} --- the latter for a peer who is present and willing but
has nothing to review right now. A peer idle past
\texttt{INACTIVITY\_WINDOW} (30 days) can be removed by anyone via the
permissionless \texttt{pruneInactivePeer} --- in practice the
\texttt{consensus-keeper} (Section~15) runs it automatically on a schedule,
but because the call is permissionless the keeper is a convenience, not a
dependency. Because inactivity is objective and on-chain, no vote is needed
(unlike the subjective misbehaviour revoke), and the prune can never drop the
active set below the seed-phase floor, so garbage collection can neither
strand nor capture the network.

The same permissionless-after-the-window discipline cleans up every other
transient state: \texttt{lapseNominee} clears a dead nomination,
\texttt{cancelStaleRevocation} clears a dead revoke motion,
\texttt{lapseProposal} clears a dead taxonomy proposal,
\texttt{cancelStaleRetire} clears a dead retire motion, and
\texttt{markLapsed} / \texttt{finalizeChallenge} settle a binding whose
window has closed. None of these require privilege; the contract never
depends on a caretaker to make progress.

% =====================================================================
\section{Taxonomy retirement}
% =====================================================================

A ratified node is never deleted --- the log is immutable --- but a strong
peer supermajority can \emph{retire} a node so it drops out of the public
taxonomy. Retirement mirrors the revocation motion:
\texttt{motionRetireNode(id, noteHash, sig)} opens it (the motioner is
retire-vote one, signing an EIP-712 \texttt{Vote} of phase~3
\texttt{retire} for the new round, with \texttt{bindingId} = the node
id), \texttt{voteRetireNode(id, noteHash, sig)} accrues further signed
votes, and at $\ceilop{2n/3}$ the node flips to \texttt{Retired} and is
removed from the public lists. A stale motion is cleared by
\texttt{cancelStaleRetire} after the window.

The mechanism enforces ``no empty pillars'' from the other direction. Only
\textbf{topics} are retired directly; \texttt{motionRetireNode} reverts on a
pillar. A pillar is retired \emph{automatically}, in the same transaction,
when its \emph{last} topic is retired --- so a pillar can never sit ratified
with zero topics. (Symmetrically, a topic proposal that would attach to a
pillar retired while the proposal was in flight is not ratified; it is left
to lapse, so a simple-majority proposal can never partially undo a $2/3$
retirement.) The off-chain indexer flips both rows to \texttt{retired} as it
handles each \texttt{NodeRetired} event.

% =====================================================================
\section{Universal vote-by-signature: every peer act is signed}
% =====================================================================

Every act a peer can take on this system --- approving or rejecting a
record, opening or voting in a challenge, proposing or endorsing a
pillar/topic, motioning or voting on a retirement, motioning or voting
on a force-renounce, nominating a peer, endorsing a nominee, or
motioning/voting to revoke a peer --- is authorised by an EIP-712
typed-data signature that the contracts recover \emph{on-chain} and
match against \texttt{msg.sender}. The signed message is not a
description of the vote that lives only off-chain; it \emph{is} the
authorisation the contract checks before mutating state. Two typehashes
carry the whole surface:

\textbf{Core \texttt{Vote} --- evidence, taxonomy, retirement,
force-renounce.} The core contract recovers signatures against a single
typehash, keyed by a \texttt{phase} field that names the act:

\begin{plainbox}
\small
\texttt{domain = \{ name: "EvidenceConsensus", version: "1",}\\
\hspace*{2em}\texttt{chainId, verifyingContract = core \}}\\[0.5em]
\texttt{Vote \{}\\
\hspace*{1.5em}\texttt{bindingId : bytes32}\\
\hspace*{1.5em}\texttt{phase     : uint8}\;\; // 0 review, 1 challenge,\\
\hspace*{1.5em}\hspace*{8.5em}    // 2 taxonomy, 3 retire, 4 force-renounce\\
\hspace*{1.5em}\texttt{support   : bool}\\
\hspace*{1.5em}\texttt{round     : uint32}\\
\hspace*{1.5em}\texttt{noteHash  : bytes32}\\
\texttt{\}}
\end{plainbox}

For phase \texttt{review} / \texttt{challenge} the \texttt{bindingId} is
the (evidence~$\times$~topic) id; for phase \texttt{taxonomy} it is the
node id (a pillar or topic propose/endorse); for phase \texttt{retire}
it is the node id of the retirement target; for phase
\texttt{force-renounce} there is no subject node, so a fixed sentinel
\texttt{keccak256("force-renounce")} stands in. The \texttt{round}
field binds the signature to the current motion / review / challenge
round (anti-replay across re-filings and re-motions), and a propose /
motion (which mints the new round) signs the round it is about to
create (\texttt{current + 1}). \texttt{noteHash} commits an optional
off-chain deliberation note (\texttt{0x00...00} for ``no note'').

\textbf{Governance \texttt{PeerVote} --- nominate, endorse-nominee,
revoke.} The membership sidecar uses its own domain and typehash so the
two signing surfaces cannot be cross-replayed:

\begin{plainbox}
\small
\texttt{domain = \{ name: "PeerGovernance", version: "1",}\\
\hspace*{2em}\texttt{chainId, verifyingContract = sidecar \}}\\[0.5em]
\texttt{PeerVote \{}\\
\hspace*{1.5em}\texttt{subject  : address}\\
\hspace*{1.5em}\texttt{kind     : uint8}\;\; // 0 nominee-endorse, 1 revoke,\\
\hspace*{1.5em}\hspace*{8em}     // 2 nominate (the nominator's signed act)\\
\hspace*{1.5em}\texttt{support  : bool}\\
\hspace*{1.5em}\texttt{round    : uint32}\\
\hspace*{1.5em}\texttt{noteHash : bytes32}\\
\texttt{\}}
\end{plainbox}

In both structs the signature is produced through the peer's wallet
(\texttt{eth\_signTypedData\_v4}) and recovered by the contract via
\texttt{ecrecover} against its own \texttt{DOMAIN\_SEPARATOR} and the
corresponding typehash. A signature that does not recover to
\texttt{msg.sender} reverts the call. There is no path by which a third
party --- the platform included --- can post a vote on a peer's behalf,
because the call would not carry that peer's signature.

\textbf{The deliberation note.} Every signed act may carry an optional
note that captures the peer's reasoning. The note \emph{text} cannot
live on-chain (it would not fit, and storage costs would discourage
detail), so what rides in the signature is the note's keccak-256
fingerprint \texttt{noteHash}. The text is then submitted to one of
four edge functions, each of which re-verifies the same signature the
chain recovered, re-hashes the submitted note and checks it matches
\texttt{noteHash}, and only then writes the row:
\texttt{verify-attestation} (review, challenge, taxonomy endorse, with
the taxonomy row also carrying \texttt{node\_hash} + \texttt{round} so
the public proof modal recovers the underlying \texttt{Vote});
\texttt{nominee-vote} (nominate + nominee-endorse, both
\texttt{PeerVote}); \texttt{revocation-vote} (revoke motion + discard,
both \texttt{PeerVote}, plus an off-chain ``keep'' dissent that has no
on-chain counterpart and is stored as a signed Attestation); and
\texttt{gov-note} (the retire + force-renounce notes, both phase~3 / 4
\texttt{Vote}s, persisted alongside the matching on-chain vote event).

\begin{notebox}
\textbf{Why every act, why on-chain.} The chain by itself can record
\emph{that} a vote happened. By recovering the signature on-chain we
make a stronger claim: the call could only have come from the wallet
that signed the typed message. The deliberation note --- committed by
hash on-chain, attached as readable text off-chain --- means a peer's
reasoning is bound to the same act they were counted for, not a
detachable comment they could later distance themselves from.
``Judgement under your own name'' becomes ``judgement under your own
\emph{signature}, with your reasons attached, recoverable by anyone.''
\end{notebox}

\textbf{Off-chain Attestations are still used as a public proof
artefact.} Alongside on-chain recovery, every counted vote also lands
in the off-chain \texttt{attestations} table as an EIP-712
\texttt{Attestation} (the same typed message used in the original
scheme: \texttt{evidenceId / topicId / peerAddr / phase / verdict /
note}). Two reasons keep this:

\begin{itemize}[leftmargin=1.4em]
  \item \textbf{Public proof, browser-verifiable.} A visitor can re-run
  the identical \texttt{Attestation} recovery in their own browser over
  the public \texttt{(domain, types, message, signature)} and confirm
  the \emph{named peer} authored the vote --- no need to ask the chain,
  no trust in the platform's UI required.
  \item \textbf{Dissents the chain cannot represent.} A taxonomy
  \emph{reject} and a revocation \emph{keep} have no on-chain
  counterpart (the contract has no reject-node or keep-peer call), so
  they are admitted as signed \texttt{Attestation} dissents with no
  tally effect --- a public record of disagreement that the contract
  itself does not need to count.
\end{itemize}

The off-chain store keeps exactly one row per \texttt{(binding\_id,
peer\_addr, phase)}, mirroring the contract's no-double-voting rule.
Vote-counting in the projection acts only on the
\texttt{review} and \texttt{challenge} phases (the binding-tally
phases); a taxonomy, retire, or force-renounce attestation never moves
a binding's tallies, and the on-chain endorse / vote gate remains the
sole consensus mechanism for those.

% =====================================================================
\section{The off-chain projection}
% =====================================================================

The Supabase layer is a single, rebuildable projection. Its main tables
mirror the on-chain concepts:

\begin{center}
\begin{tabular}{@{}p{0.33\linewidth}p{0.6\linewidth}@{}}
\toprule
\textbf{Table} & \textbf{Holds} \\
\midrule
\texttt{evidence}    & Canonical content + \texttt{content\_hash}. \\
\texttt{bindings}    & One row per (evidence, topic): status, tallies, \texttt{binding\_hash}. \\
\texttt{pillars} / \texttt{topics} & Taxonomy nodes (proposed~\arr~ratified~\arr~retired), joined by \texttt{node\_hash}. \\
\texttt{attestations}& One per \texttt{(binding\_id, peer\_addr, phase)} --- the public proof artefact. \\
\texttt{nominee\_votes}    & One per (nominee, round, voter): nominate + endorse notes, \texttt{PeerVote} sig. \\
\texttt{revocation\_votes} & One per (peer, round, voter, verdict): keep + discard notes, \texttt{PeerVote} sig (keep is a signed off-chain dissent). \\
\texttt{chain\_events} / \texttt{chain\_event\_cursor} & Indexer input and progress. \\
\texttt{tamper\_alerts} & Evidence content-hash mismatches \emph{and} taxonomy node\_hash / meta\_hash drift. \\
\bottomrule
\end{tabular}
\end{center}

\vspace{0.4em}

Seven edge functions operate the layer.
\texttt{chain-indexer-evidence} reads chain events from both the core
and the governance sidecar, reconciling the projection by
\texttt{binding\_hash} and \texttt{node\_hash} (and the sidecar's
nominee/revocation events by address + round).
\texttt{verify-attestation} verifies the review/challenge/taxonomy-endorse
signature on write, re-recovers the matching on-chain \texttt{Vote},
and (for review/challenge) recomputes the content hash before recording
the attestation. \texttt{nominee-vote} verifies a nominate or
nominee-endorse \texttt{PeerVote}, matches it against the on-chain
sidecar event, and writes the note. \texttt{revocation-vote} does the
same for revoke motion / discard, and also accepts a signed
\emph{keep} dissent that has no on-chain counterpart.
\texttt{gov-note} verifies the retire + force-renounce \texttt{Vote}
and persists the note alongside the chain-emitted vote event.
\texttt{audit-content-hash} re-hashes the archive \emph{and} the
taxonomy on a schedule to detect drift --- it recomputes each evidence
\texttt{content\_hash} and each pillar/topic \texttt{node\_hash} +
\texttt{meta\_hash}, opening a \texttt{tamper\_alert} on any mismatch
and auto-resolving one that matches again. \texttt{consensus-keeper}
is described below.

Vote counts on bindings are applied through atomic database functions,
\texttt{apply\_review\_counts} and \texttt{apply\_challenge\_counts},
so concurrent votes cannot corrupt a tally. Write endpoints are
throttled by prefix in an \texttt{edge\_rate\_limit} table
(\texttt{ev\_insert:} for evidence, \texttt{tax\_insert:} for taxonomy,
\texttt{nominee:} for nominate / endorse, \texttt{revoke:} for revoke
motion / vote). Because the schema seeds nothing, the entire taxonomy
and archive are produced at runtime by peers, never shipped as a
baseline.

\textbf{The keeper.} \texttt{consensus-keeper} is a \texttt{pg\_cron}-scheduled
edge function (every five minutes) that drives the two on-chain maintenance
jobs the contract cannot trigger for itself: it prunes peers idle past
\texttt{INACTIVITY\_WINDOW} (\texttt{pruneInactivePeer}, oldest-idle first,
stopping at the seed-phase floor) and promotes queued bindings into review in
public-priority order (\texttt{promote}, until \texttt{reviewCapacity} is
full). It holds a signing key (\texttt{KEEPER\_PRIVATE\_KEY}) to send these
transactions, but it grants \emph{no} authority: both calls are permissionless
and objectively gated on-chain, so the keeper can do nothing any address could
not already do, and anyone may run the same jobs if it goes offline. It is a
convenience that automates liveness, never a privileged actor.

% =====================================================================
\section{A note on residual risk: open-submission front-running}
% =====================================================================

Honesty about limits is part of the design. \texttt{submitEvidence} is
permissionless, and on a public chain an evidence id (the Supabase UUID) is
visible in the mempool before it confirms. An adversary can therefore
front-run a submission with the \emph{same} id but a different content hash,
reverting the victim (``already submitted'') and binding the id to junk
content. This is bounded, not catastrophic:

\begin{itemize}[leftmargin=1.4em]
  \item the off-chain \texttt{audit-content-hash} job flags any on-chain hash
  that does not match the canonical payload, so the junk is detected;
  \item UUIDs are free, so the victim simply re-mints and resubmits;
  \item nothing enters the public archive without passing peer review, and
  taxonomy proposals are additionally peer-gated --- any squatted-but-
  unratifiable node is freed by the permissionless \texttt{lapseProposal}
  after the window.
\end{itemize}

A commit-reveal submission scheme would close the window entirely, at the
cost of a second transaction per submission; it is deliberately omitted here.
Documenting the residual rather than hiding it is the same instinct that
animates the whole system: the gate is a published rule, and so are its
edges.

% =====================================================================
\section{How the engineering serves the philosophy}
% =====================================================================

Every mechanism above exists to make a stated principle enforceable rather
than merely promised.

\begin{center}
\begin{tabular}{@{}p{0.46\linewidth}p{0.46\linewidth}@{}}
\toprule
\textbf{Principle (the essay)} & \textbf{Mechanism (this paper)} \\
\midrule
Filed against a source &
Content-only \texttt{contentHash}, recomputed in three places; mismatches
raise tamper alerts. \\[0.4em]
Judged in public by named peers &
Every peer act is an EIP-712 \texttt{Vote} or \texttt{PeerVote} recovered
\emph{on-chain} from the voter's wallet signature; re-verifiable by
anyone, no anonymous moderators, no third-party-posted votes. \\[0.4em]
Reasons attached, not just verdicts &
Each signed act commits a \texttt{noteHash}; the note text is recorded
off-chain by four edge functions that re-verify the same signature, so
every counted vote carries its peer's reasoning. \\[0.4em]
Canon is a majority, not a faction &
True-majority canonize thresholds (60 / 55 / 51\%); order-independent via
peer-count snapshots. \\[0.4em]
Democratized method, open intake &
Permissionless \texttt{submitEvidence} with a fixed-batch review queue;
an administrator-free registry on the \texttt{PeerGovernance} sidecar
that admits and removes peers by signed consensus. \\[0.4em]
Evidence is not proof &
Tiered evidence with asymmetric thresholds; ``Canon'' is a quorum
judgement, never a truth flag. \\[0.4em]
Evolving consensus reality &
Peer-grown Pillar~\arr~Topic taxonomy; a perpetual challenge route;
reaffirm-or-deprecate revision; supermajority retirement. \\[0.4em]
Owned by no one, seizable by no one &
An open core + Lens + peer-gated sidecar as the sole authority; two-step
transfer, renounce, and a $2/3$ \emph{signed} force-renounce that evicts
a captured owner. \\[0.4em]
Capture-resistant &
A strict $1/3{+}1$ admission gate (BFT honest-supermajority boundary)
and a liveness clock that keeps the quorum denominator honest. \\
\bottomrule
\end{tabular}
\end{center}

\vspace{0.6em}

The contracts are the place where the philosophy stops being a
description of intentions and becomes a description of what the system
can and cannot do. ``Anyone may file'' is a function anyone can call.
``Named peers verify'' is an EIP-712 signature the contract itself
recovers, not a database row a moderator could fabricate. ``With their
reasons'' is a \texttt{noteHash} riding in that signature. ``A majority
stands behind it'' is a counted vote against a frozen denominator.
``Owned by no one'' is a key that two thirds of peers can take from any
hand. ``Revised forever'' is a state machine with no terminal lock on
canon. The mechanism is the argument that the values are meant --- and
the companion essay, \emph{Disclosure: A Labour of Love}, is the
argument for why they are worth meaning in the first place.

\begin{notebox}
\textbf{Open source.} None of this has to be taken on trust --- every claim in
this paper is checkable against the code. The smart contract, its Lens
sidecar, the edge functions, the frontend, and the sources of both
whitepapers live in the public repository:\\[0.4em]
\centering\url{https://github.com/DisclosureG/the-disclosure-platform}
\end{notebox}

\end{document}
